\chapter{RG-инвариантное чтение дискретного разрешения}
\label{ch:v10-discrete-resolution-rg-readout}

Нормировки $\ell_{\rm cell}=1$ воксель и
$\tau_{\rm birth}=1$ такт понимаются как внутренние единицы, а не как
утверждение о числе метров. Наблюдаемой величиной служит разрешение моды
\begin{equation}
 \mathcal R_\mu=\frac1{\mu_{\rm spec}\ell_{\rm cell}}.
\end{equation}
Два предыдущих точных результата дают
\begin{equation}
 \Lambda_{43}\ell_{\rm cell}=42,
 \qquad
 \frac{\mu_{\rm spec}}{\Lambda_{43}}
 =\frac{e^{-32\pi^2/3}}{42}.
\end{equation}
Поэтому число 42 сокращается:
\begin{equation}
 \mu_{\rm spec}\ell_{\rm cell}=e^{-32\pi^2/3},
 \qquad
 \boxed{\mathcal R_\mu=e^{32\pi^2/3}}.
\end{equation}

При масштабной орбите
$\ell_{\rm cell}\mapsto s\ell_{\rm cell}$,
$\mu_{\rm spec}\mapsto\mu_{\rm spec}/s$ величина $\mathcal R_\mu$
не меняется. Если $\ell_{\rm cell}=c\tau_{\rm birth}$, то время моды
$\tau_\mu=1/(c\mu_{\rm spec})$ удовлетворяет
\begin{equation}
 \frac{\tau_\mu}{\tau_{\rm birth}}=\mathcal R_\mu.
\end{equation}
Таким образом пространственное и временное разрешения совпадают.

Логарифмическая карта двух исходных отношений имеет ранг/ядро $2/1$.
Строка для $\mu_{\rm spec}\ell_{\rm cell}$ является их суммой, поэтому
добавление разрешения не повышает ранг и не определяет абсолютный размер
вокселя. Оно зато даёт настоящий инвариант масштабной орбиты.

Численно $\mathcal R_\mu\approx5.2564\times10^{45}$. Следовательно,
отождествления с протонным $10^{20}$ или секундным $10^{43}$ пока не
выведены: это безразмерные несовпадения, которые нельзя убрать сменой единиц.

\begin{theorem}[Принцип дискретного разрешения]
При условной типизации $\mu_{\rm spec}$ как обратной комптоновской длины
физическое содержание масштабной ветви выражается не абсолютным
$\ell_{\rm cell}$, а инвариантным числом клеток и тактов
$\mathcal R_\mu=e^{32\pi^2/3}$.
\end{theorem}

\begin{nogocriterion}[Единица вокселя не есть метрический вывод]
Равенство $\ell_{\rm cell}=1$ фиксирует внутреннюю систему единиц. Оно не
доказывает планковскую длину и не разрешает приписывать моде имя протона без
независимого морфизма $\mu_{\rm spec}=mc/\hbar$ для конкретного состояния.
\end{nogocriterion}

\begin{protocol}\begin{enumerate}
\item условная архитектура чтения: \textbf{$10/10$};
\item точные реляционные следствия: \textbf{$7/7$};
\item сокращение K43-фактора: \textbf{$1/1$};
\item масштабная инвариантность: \textbf{$1/1$};
\item типизация физической моды: \textbf{$0/3$};
\item необходимость абсолютной линейки для $\mathcal R$: \textbf{нет};
\item следующий гейт: \textbf{аудит физической моды разрешения}.
\end{enumerate}\end{protocol}

Машинный сертификат сохранён в
\path{s2t_v10_cell_birth_four_volume_nonequilibrium_bath_discrete_resolution_rg_invariant_readout_gate_results.json}.